\documentclass{IOS-Book-Article}
\usepackage{mathptmx}
\usepackage{multirow}
\usepackage{graphicx}
\usepackage[utf8]{inputenc}
\usepackage[T1]{fontenc}
\usepackage[english]{babel}
\begin{document}
\begin{frontmatter}
\title{Synthetic corpora in the adaptation of language models to Brazilian law}
\author[A]{\fnms{Eduardo} \snm{Pacheco}}
and
\author[A]{\fnms{Cibele Maria} \snm{Russo}}
\runningauthor{E. Pacheco and C. M. Russo}
\address[A]{Instituto de Ciências Matemáticas e de Computação, Universidade de São Paulo}
\begin{abstract}
The standard recipe for adapting a language model to a technical domain has two stages: expose the model to domain text under next-token prediction, then to examples of the solved task. The recipe does not specify \emph{which} text, \emph{in what form} it should enter, nor \emph{how much} supervision suffices. This work measures the three at once, in a domain where an itemised and public grading rubric exists: the second phase of the Brazilian Bar Examination, the entrance examination for the legal profession in Brazil. We built two continued pre-training corpora paired document by document, one of raw legal text and one in which the same material was restructured into cards with layers generated by a teacher model. We crossed these two regimes, plus a control regime without pre-training, with six fine-tuning sets of distinct composition and with four base models from 1.5 to 14 billion parameters, for a total of 180 conditions trained with LoRA. Each condition was evaluated on 105 real questions against 940 criteria from the examining boards' official rubric, over ten independent runs, judged by a language model calibrated against human annotation. We obtain three results. First, the effect of each decision depends on the textual genre of the task and cancels out when genres are aggregated: continued pre-training yields between +2.6 and +3.9 percentage points on legal document writing and essentially nothing on discursive questions. Second, the two corpus regimes are \textbf{not substitutes}: raw text teaches document form and the structured corpus teaches answering questions, with a correlation of only $\rho$ = 0.30 between their criterion-level gains. Third, pre-training on raw text \textbf{degrades} the model before fine-tuning, by 9.8 percentage points on average in document writing, so that the observed gain is deficit recovery followed by overtaking, and not an increment over a tie.
\end{abstract}
\begin{keyword}
domain adaptation \sep Brazilian law \sep LLM-as-a-judge evaluation \sep LoRA
\end{keyword}
\end{frontmatter}

\section{Introduction}
Adapting a general-purpose language model to a technical domain follows a well-known recipe. One takes a pre-trained model, exposes it to domain text under next-token prediction, then to examples of the solved task, and measures the outcome. The literature repeats this recipe in medicine, finance and law, and the reports converge on the same qualitative conclusion: adaptation helps.
What the recipe usually does not say is \textbf{which text}. In Brazilian law there are at least four distinct categories of legal text that could compose a training corpus, and they differ considerably. Statutory text is normative and compact, enumerating norms without explanatory intent. Case law is argumentative and repetitive. Doctrine is explanatory. Professional entrance examinations require candidates to write texts matching the skills expected of legal practitioners. It is not obvious that a model learns the same thing from each category, nor that it needs all of them.
The second thing the recipe does not say is \textbf{in what form}. A statutory article may enter training exactly as published, in raw form. Or it may be reshaped so as to maximise its association with legally relevant concepts, forming a kind of "study card" of the sort law students write. Instead of appearing sequentially alongside the other articles of the same statute, the article gets its own focus and comes with an explanation, a summary, its practical effect, its relation to other articles, and question-answer pairs derived from it. The second form has a cost: the material must be generated with another language model, and it takes more room in the token budget. On one hand, enumerating statutory articles in raw form admits more articles and more statutes, given their compact format. On the other, spending effort on establishing these connections may benefit model capabilities on legal tasks. The question is whether, and under which circumstances, each format pays off.
The third question is whether the answers to the first two \textbf{depend on scale}. A result obtained on a 1.5-billion-parameter model is only useful if it still holds at 14 billion, and there is theoretical reason to doubt that it does: a larger model has already seen more legal Portuguese in its original pre-training, so additional material may be redundant for it and novel for the smaller one.
This work answers the three questions with a factorial design. The contribution is not to show that adaptation helps, which is already known, but to decompose where the gain resides. The second phase of the Bar Examination assesses two tasks, (a) answering discursive questions and (b) drafting a legal document, and they respond in opposite directions to nearly every decision tested. A score that sums the two measures the difference between a positive and a negative effect, and thus the effect of nothing at all.
The two pre-training corpora and the six fine-tuning blocks [11], the judgments of the 180 conditions together with the responses they generated [12], and the code [14] are openly deposited, with per-file checksums and a datasheet describing the collection and the known defects.
\section{Related work}
Continuing the pre-training of a model on domain text before fine-tuning was systematised by Gururangan et al. [1], who showed consistent gains across four domains. Later work applied the procedure to law: Chalkidis et al. [2] trained encoders on English-language legal corpora and Colombo et al. [3] extended the recipe to generative models. In all these cases the domain corpus enters in the form in which it was collected, and the question of \emph{restructuring} it before training is not posed.
For legal Portuguese there are three directly relevant evaluation sets, all built on real examinations, and two of them on the same one we use. \textbf{oab-bench} [4] draws 105 questions across seven areas of law from recent editions of the Bar Examination and grades them on a 0 to 10 scale against the examiners' guidelines, reporting a strong correlation between a frontier model's scores and human ones. \textbf{Magis-Bench} [5] moves to the judiciary, with 74 questions from eight competitive examinations for judicial positions between 2023 and 2025, covering discursive analysis and the drafting of civil and criminal judgments, graded by four frontier models acting as judges. \textbf{Rabula} [6] is built on the 2024 Bar Examination with four tasks, and grades the generative ones against the board's own criteria as binary questions weighted by difficulty, calibrating the automatic judge by Cohen's agreement with three legal experts.
This work inherits its evaluation set and grading pipeline from [6] and does not modify them. The distinction from [4] is worth stating, since both rest on the same examination: [4] scores a response as a whole from 0 to 10, whereas the pipeline used here decides each of the 940 rubric criteria separately, which is what allows performance to be reported by criterion class rather than by question. The results below depend on that decomposition.
Evaluating long-form answers with a language model [7] accommodates the format of the examination at the cost of raising the question of how far the judge agrees with the human expert. We address that question in a companion paper [8]; here the judges enter already fixed and calibrated, and are not an object of study. As for decomposing legal document writing into requirements of form and requirements of merit, it does not appear in the domain-adaptation literature we surveyed, and it is from that decomposition that most of the effects measured below come.
\section{Materials and methods}
Four base models enter the experiment, in two families and four scales: Tucano2 at 1.5 and 3.7 billion parameters, trained with a focus on Portuguese, and Qwen3 at 8.2 and 14.8 billion, multilingual. The family difference is deliberate, because it allows a scale effect to be distinguished from an effect of prior exposure to legal Portuguese. One fine-tuning condition costs between 4.8 and 28 hours on an NVIDIA A10G GPU depending on size.
\subsection{The pre-training corpus: the raw and synthetic-card pair}
The central axis of the design is the form of the pre-training corpus. We built two corpora from the same collection of documents gathered from official sources: federal and state statutes, appellate decisions and summary opinions from the higher courts and from the São Paulo State Court of Appeals. The two are paired document by document: the same decision that enters one enters the other.
In the \textbf{raw text} regime the document enters as collected and is cut by length when it exceeds the context window. The cut is blind: it falls wherever it falls, including mid-sentence. There are 39,856 documents and 50.58 million tokens under the Tucano2 tokeniser. The composition is dominated by judicial decisions: 78.3\% of tokens are full appellate decisions, 9.1\% consolidated case law (summary opinions and binding precedent theses), 8.9\% statutory text and 3.7\% enacted opinions.
In the \textbf{structured} regime the same document is converted into a \textbf{card} before entering. The card preserves the original text intact in a field of its own and adds six layers generated by a teacher model: a summary, an explanation, the practical effect, related provisions, the principles involved, and question-answer pairs derived from the content. Each layer is a node with an identifier and an explicit parent relation, which allows the card to be sliced without losing its anchor: the cut respects node boundaries, and each piece carries the identification header plus a coherent subset of the layers. There are 41,900 documents and 59.43 million tokens.
Expansion inverts the composition, and this is a limitation of the pair. A statute has few tokens in the original and many once expanded; a full appellate decision is already long and expansion adds proportionally less. The structured arm ends up with 36.3\% statute and 35.5\% case decision, against 8.9\% and 78.3\% in the raw arm. The two start from the same collection but do not reach the GPU with the same mixture of categories, so part of the effect attributed to form may be mixture.
The token budget is that of the smaller arm, counted with each family's tokeniser. The raw arm runs exactly one epoch over the whole file; the structured arm is truncated to match the same token count, which gives 0.851 epochs on Tucano2 and 0.877 on Qwen3. Neither arm repeats a document.
\subsection{The fine-tuning corpus: what each block contains}
The fine-tuning corpus consists of prompt-response pairs, and the examples vary along two independent axes. The first is the origin of the prompt: it may be a real question that appeared in the examination, or a synthetic question written by a model from a real case used as a seed and, in the extreme, a synthetic question written from another synthetic question, two hops away from official material. The second is the anchor of the response: it may follow the board's official rubric, with the published criteria and points, or be free, with a rubric generated alongside it by the teacher model.
Five sets account for 2,630 of the 2,909 examples and give the blocks their identity:
\begin{itemize}
\item \textbf{D1} \textperiodcentered{} 691 real \textbf{discursive questions} from the second phase of the Bar Examination, with the official rubric as the response. Prompt at zero hops, official anchor. Median response of 376 tokens.
\item \textbf{D5} \textperiodcentered{} 181 real \textbf{legal documents} from the second phase, also with the official rubric. Median response of 1,579 tokens, the longest set in the corpus.
\item \textbf{D3} \textperiodcentered{} 558 questions from the \textbf{first phase} of the Bar Examination, multiple-choice, converted into open questions. Prompt at one hop, and the anchor is the objective answer key, which states the correct alternative but does not develop the reasoning.
\item \textbf{D2} \textperiodcentered{} 686 \textbf{synthetic legal documents}, written from a real second-phase case used as a seed. Prompt at one hop, free anchor.
\item \textbf{D4} \textperiodcentered{} 514 \textbf{cascaded synthetic documents}, written from a D3 prompt. Prompt at two hops, free anchor.
\end{itemize}
The remaining 279 examples come from nineteen small sets drawn from \textbf{other legal career examinations}: documents, legal opinions, judgments and essays from various boards, from the Office of the Attorney General and from the São Paulo judiciary.
The six blocks combine these sets, and they do not form a ladder of volume:
\begin{itemize}
\item \textbf{I} \textperiodcentered{} 1,007 examples, official Bar material with the board's rubric only (D1 and D5).
\item \textbf{II} \textperiodcentered{} 1,182, generated or converted material only (D2 and D3). It does not contain block I.
\item \textbf{III} \textperiodcentered{} 1,681, block II plus the two-hop material (D4).
\item \textbf{IV} \textperiodcentered{} 2,687, block I plus block III, that is, all Bar material, official and generated.
\item \textbf{V} \textperiodcentered{} 223, other career examinations only, no Bar material.
\item \textbf{VI} \textperiodcentered{} 2,909, block IV plus block V: the complete corpus.
\end{itemize}
Nesting was verified by content, not by label: II $\subset$ III $\subset$ IV $\subset$ VI and I $\subset$ IV, but \textbf{I and II are disjoint}, as are I and III.
\subsection{Factorial design and the three starting points}
The experiment crosses three factors: the pre-training regime, at three levels, (i) raw text, (ii) structured and (iii) none, which skips the stage entirely and serves as control; the fine-tuning block, at six levels; and the base model, at four levels. Each combination was trained with three distinct random seeds, for a total of 180 conditions.
The design has three untreated starting points, and each answers a different question. The \textbf{base model} is the model as published, without pre-training and without fine-tuning; comparing it with a trained condition measures the whole pipeline. \textbf{Block 0} is the model after continued pre-training and before any task example, and it exists in two versions, one per corpus regime; comparing it with a fine-tuning block measures fine-tuning. The \textbf{none} regime is the model that skipped pre-training and went straight to fine-tuning; comparing it with the same block under the other regimes measures pre-training. Note that the \texttt{none} regime is not the same model across blocks, because it removes pre-training and not fine-tuning; the only starting point that is constant across blocks is block 0.
All 180 conditions use the same protocol: low-rank adaptation (LoRA) [9] with rank 32 and scaling factor 64, applied to the last 24 layers and to seven projection modules; learning rate of 1$\times$10$^{-}$$^{5}$ with cosine decay; effective batch of 16; a 3,072-token window; and prompt masking, so that the loss is computed over the response only. The number of steps is not fixed: it is computed so that every block receives exactly three epochs, which makes the number of examples, and not the number of updates, the manipulated variable.
\subsection{Evaluation}
Evaluation follows the methodology published in earlier work [6] and rests on \emph{LLM-as-a-judge} with models that rank well on legal document writing and on answering discursive questions, comparing several automatic judges against human judgment on whether a model's responses satisfy the official rubrics of the second phase of the Bar Examination.
The evaluation set is fixed and was never touched by training: \textbf{105 questions} from the second phase, 84 discursive and 21 practical cases requiring a legal document. The examinations they come from were excluded from the entire training corpus, which we verified by examination identifier in each block.
Each question is graded against the board's official rubric, item by item. There are \textbf{940 binary criteria} in total, each with the points the board assigns to it. The score of a condition is the per-run average of the fraction of points obtained on each question, rather than the plain proportion of criteria satisfied, which would give more weight to questions with longer rubrics. Each condition is generated ten times independently, at temperature 0.1.
There are two judges, one per genre, chosen and fixed before any measurement of the 180 conditions: \textbf{Kimi K2.5} for discursive questions, with Cohen's agreement of \textbf{$\kappa$ = 0.799} against the human gold standard, and \textbf{GPT-OSS 120B} for document writing, with \textbf{$\kappa$ = 0.754}. For reference, agreement among the three human annotators is 0.789 on discursive questions and 0.763 on document writing.
The board's rubric gives each document-writing criterion a title: Addressing, Party identification, Closing, Legal grounds, Requests. We grouped these titles into two classes: those of \textbf{form}, which demand the skeleton of the document and amount to 148 criteria and 14.9\% of the points, and those of \textbf{merit}, which demand the legal argument and amount to 413 criteria and 85.1\%. The titles are the board's; the grouping is ours. We report formal, merit and aggregate document scores whenever document writing is measured.
Conditions are compared by Wilcoxon signed-rank tests paired by question, with the three seeds stacked. Where the number of comparisons is large, Bonferroni and Benjamini-Hochberg corrections are applied [10].
\section{Results}
The tables below give the score of every condition, as a percentage of the points the board distributes. Each value is the mean of three seeds, and each seed is the mean of ten runs. The \textbf{block 0} column is the model before any fine-tuning, and the \textbf{no pre-training} row is the control: these are the two zeros against which the rest is read. The best block of each row is highlighted. The blocks have I = 1,007 examples, II = 1,182, III = 1,681, IV = 2,687, V = 223 and VI = 2,909. The \emph{no pre-training} row was run across all six blocks only on Tucano2 3.7B; on the other three models it exists for blocks I and VI.
As for dispersion, the standard deviation across the ten runs of a seed lies between 1.8 points on discursive questions and 3.0 on document merit; since each cell averages three seeds, what matters for comparing conditions is the dispersion between seeds, from 0.4 to 0.9. The execution noise of the protocol, measured over 340 accumulated runs, is 1.20 points: smaller differences are not distinguishable from fluctuation.
\begin{table}[!ht]
\centering
\caption{Model performance on discursive questions.}
\footnotesize\setlength{\tabcolsep}{3pt}
\begin{tabular}{llrrrrrrr}
\hline
\textbf{Model} & \textbf{Pre-training} & \textbf{block 0} & \textbf{I} & \textbf{II} & \textbf{III} & \textbf{IV} & \textbf{V} & \textbf{VI} \\
\hline
\multirow{3}{*}{Tucano 1.5B} & no pre-training & 13.8\% & 14.7\% & -- & -- & -- & -- & 14.3\% \\
 & raw text & 11.1\% & 15.5\% & 15.7\% & 16.3\% & 14.9\% & 13.4\% & 13.9\% \\
 & cards & 17.7\% & 15.8\% & 16.5\% & 17.7\% & 15.2\% & 15.5\% & 16.0\% \\
\multirow{3}{*}{Tucano 3.7B} & no pre-training & 25.0\% & 26.8\% & 29.5\% & 29.4\% & 27.7\% & 26.5\% & 28.3\% \\
 & raw text & 22.1\% & 26.5\% & 29.9\% & 29.5\% & 28.4\% & 26.4\% & 27.8\% \\
 & cards & 31.0\% & 28.1\% & 31.2\% & 31.8\% & 29.1\% & 29.7\% & 28.6\% \\
\multirow{3}{*}{Qwen3 8B} & no pre-training & 29.5\% & 26.8\% & -- & -- & -- & -- & 27.3\% \\
 & raw text & 28.2\% & 24.5\% & 26.2\% & 24.8\% & 25.1\% & 29.4\% & 24.2\% \\
 & cards & 30.8\% & 26.4\% & 26.2\% & 26.2\% & 26.1\% & 30.4\% & 25.7\% \\
\multirow{3}{*}{Qwen3 14B} & no pre-training & 35.5\% & 28.9\% & -- & -- & -- & -- & 29.7\% \\
 & raw text & 36.0\% & 31.2\% & 28.2\% & 28.3\% & 30.7\% & 26.7\% & 30.2\% \\
 & cards & 35.6\% & 31.2\% & 31.0\% & 30.8\% & 31.6\% & 28.7\% & 30.8\% \\
\hline
\end{tabular}
\end{table}

\begin{table}[!ht]
\centering
\caption{Model performance on legal document writing, formal criteria.}
\footnotesize\setlength{\tabcolsep}{3pt}
\begin{tabular}{llrrrrrrr}
\hline
\textbf{Model} & \textbf{Pre-training} & \textbf{block 0} & \textbf{I} & \textbf{II} & \textbf{III} & \textbf{IV} & \textbf{V} & \textbf{VI} \\
\hline
\multirow{3}{*}{Tucano 1.5B} & no pre-training & 14.6\% & 14.2\% & -- & -- & -- & -- & 31.1\% \\
 & raw text & 3.0\% & 26.9\% & 26.9\% & 32.3\% & 36.4\% & 8.7\% & 37.5\% \\
 & cards & 6.2\% & 23.0\% & 19.9\% & 24.4\% & 26.5\% & 9.5\% & 27.1\% \\
\multirow{3}{*}{Tucano 3.7B} & no pre-training & 21.9\% & 27.3\% & 33.3\% & 38.6\% & 39.7\% & 14.7\% & 39.2\% \\
 & raw text & 6.2\% & 36.0\% & 35.6\% & 41.3\% & 40.4\% & 21.5\% & 41.1\% \\
 & cards & 14.7\% & 34.6\% & 36.8\% & 40.2\% & 39.9\% & 16.4\% & 40.4\% \\
\multirow{3}{*}{Qwen3 8B} & no pre-training & 16.4\% & 9.2\% & -- & -- & -- & -- & 28.8\% \\
 & raw text & 9.9\% & 30.6\% & 28.0\% & 31.1\% & 33.1\% & 14.0\% & 34.6\% \\
 & cards & 18.1\% & 26.9\% & 28.6\% & 31.1\% & 31.5\% & 16.2\% & 32.7\% \\
\multirow{3}{*}{Qwen3 14B} & no pre-training & 21.4\% & 22.8\% & -- & -- & -- & -- & 31.6\% \\
 & raw text & 17.5\% & 33.4\% & 31.1\% & 35.7\% & 38.9\% & 24.8\% & 37.4\% \\
 & cards & 22.7\% & 31.7\% & 29.8\% & 35.2\% & 37.0\% & 20.4\% & 36.8\% \\
\hline
\end{tabular}
\end{table}

\begin{table}[!ht]
\centering
\caption{Model performance on legal document writing, merit criteria.}
\footnotesize\setlength{\tabcolsep}{3pt}
\begin{tabular}{llrrrrrrr}
\hline
\textbf{Model} & \textbf{Pre-training} & \textbf{block 0} & \textbf{I} & \textbf{II} & \textbf{III} & \textbf{IV} & \textbf{V} & \textbf{VI} \\
\hline
\multirow{3}{*}{Tucano 1.5B} & no pre-training & 13.5\% & 12.8\% & -- & -- & -- & -- & 15.4\% \\
 & raw text & 2.5\% & 12.2\% & 13.7\% & 15.7\% & 17.5\% & 8.4\% & 19.1\% \\
 & cards & 10.9\% & 13.1\% & 9.7\% & 12.8\% & 14.9\% & 8.7\% & 15.6\% \\
\multirow{3}{*}{Tucano 3.7B} & no pre-training & 23.3\% & 16.7\% & 26.0\% & 26.2\% & 27.7\% & 15.3\% & 27.3\% \\
 & raw text & 11.0\% & 24.3\% & 25.8\% & 30.2\% & 29.0\% & 19.1\% & 28.8\% \\
 & cards & 16.5\% & 21.9\% & 26.3\% & 27.9\% & 29.1\% & 18.1\% & 29.7\% \\
\multirow{3}{*}{Qwen3 8B} & no pre-training & 17.0\% & 4.9\% & -- & -- & -- & -- & 19.9\% \\
 & raw text & 10.5\% & 18.3\% & 19.3\% & 21.4\% & 21.5\% & 15.0\% & 22.8\% \\
 & cards & 17.6\% & 18.8\% & 21.7\% & 24.0\% & 23.2\% & 12.1\% & 24.2\% \\
\multirow{3}{*}{Qwen3 14B} & no pre-training & 18.5\% & 15.4\% & -- & -- & -- & -- & 23.4\% \\
 & raw text & 8.1\% & 22.4\% & 24.2\% & 26.3\% & 27.0\% & 20.5\% & 25.9\% \\
 & cards & 17.9\% & 24.6\% & 24.0\% & 30.3\% & 29.5\% & 18.1\% & 30.0\% \\
\hline
\end{tabular}
\end{table}

\begin{table}[!ht]
\centering
\caption{Performance by area of law on Tucano2 3.7B, averaged over the six blocks.}
\footnotesize\setlength{\tabcolsep}{3pt}
\begin{tabular}{llrrrrrrr}
\hline
\textbf{Task} & \textbf{Pre-training} & \textbf{Admin.} & \textbf{Civil} & \textbf{Const.} & \textbf{Corp.} & \textbf{Crim.} & \textbf{Labour} & \textbf{Tax} \\
\hline
\multirow{3}{*}{discursive questions} & no pre-training & 36.3\% & 23.0\% & 27.1\% & 30.4\% & 25.2\% & 21.2\% & 33.0\% \\
 & raw text & 38.7\% & 21.2\% & 24.8\% & 31.9\% & 29.3\% & 18.8\% & 31.7\% \\
 & cards & 37.8\% & 24.6\% & 29.7\% & 34.0\% & 29.0\% & 20.5\% & 32.7\% \\
\multirow{3}{*}{document, formal criteria} & no pre-training & 22.3\% & 43.5\% & 40.3\% & 45.3\% & 15.8\% & 14.3\% & 43.5\% \\
 & raw text & 29.8\% & 45.5\% & 45.0\% & 52.5\% & 16.0\% & 14.5\% & 48.5\% \\
 & cards & 26.0\% & 43.7\% & 46.7\% & 47.1\% & 21.9\% & 13.2\% & 44.3\% \\
\multirow{3}{*}{document, merit criteria} & no pre-training & 13.0\% & 24.0\% & 24.4\% & 26.4\% & 14.0\% & 22.9\% & 37.9\% \\
 & raw text & 14.3\% & 30.1\% & 25.8\% & 27.0\% & 19.1\% & 24.4\% & 42.7\% \\
 & cards & 15.8\% & 25.8\% & 33.4\% & 25.4\% & 16.6\% & 22.3\% & 39.3\% \\
\hline
\end{tabular}
\end{table}

\subsection{What each pre-training regime yields, and on which task}
Compare the \emph{no pre-training} row with the two that follow it, in the same model and the same column. On Tucano2 3.7B, averaged over the six blocks, \textbf{discursive questions} go from 28.0\% in the control to 28.1\% with raw text and 29.7\% with cards. The \textbf{formal part of document writing} goes from 32.1\% to 36.0\% with raw text and 34.7\% with cards. The \textbf{merit part} goes from 23.2\% to 26.2\% and 25.5\%.
Put differently: both pre-training regimes improve legal document writing, and neither improves discursive questions appreciably. Raw text yields +3.9 points on the formal part and +0.02 on discursive questions, which on discursive questions amounts to nothing. Cards yield +2.6 on the formal part and +1.7 on discursive questions. The pattern repeats on the other three models where a control exists. The effect size depends on the volume of supervision that follows. On block I, with 1,007 examples, raw-text pre-training takes the formal part of Tucano2 3.7B from 27.3\% to 36.0\%, a gain of 8.8 points. On block VI, with 2,909, it takes it from 39.2\% to 41.1\%, a gain of 1.9. The pre-training gain shrinks as fine-tuning grows.
The \textbf{block 0} column shows that this gain is the recovery of a deficit, and not an increment over a tie. Before any task example, pre-training on \textbf{raw text} leaves the model worse than it was: on Tucano2 3.7B the formal part falls from 21.9\% to 6.2\% and the merit part from 23.3\% to 11.0\%, and the drop repeats across all four models, averaging 9.4 points on form and 10.1 on merit. Pre-training on \textbf{cards} costs far less, 3.1 and 2.4 points, and is the only one of the two that improves discursive questions at this stage, by 2.8 points on average. A model delivered in the block-0 state with raw text would be worse than the model with no treatment at all; it is fine-tuning that recovers the deficit and, from there, overtakes the control.
\subsection{The two corpus regimes are not substitutes}
Now compare the last two sub-rows of each model. On discursive questions the cards beat raw text in nearly every column of all four models; on the formal part of document writing the reverse holds, and raw text wins in practically all of them. On the two Qwen3 models the two document cuts have \textbf{opposite signs}: cards win on merit and lose on form, and the aggregate document score, being 85.1\% merit, records only the merit sign.
The most direct explanation is compositional: a full appellate decision \emph{is} a procedural document and carries addressing, party identification and closing, whereas a card is an excerpt of a provision accompanied by questions and answers. Each regime teaches what its material contains. The division of labour was quantified at the criterion level on Tucano2 3.7B. The correlation between one regime's gain and the other's, criterion by criterion, is \textbf{$\rho$ = +0.30} (Spearman; Pearson r = +0.37), closer to zero than to one. Of the criteria improved by at least one regime, \textbf{218 are exclusive to one of them against 103 shared}. Taking, criterion by criterion, the better of the two rates yields an upper bound for the union that sits 4.92 points above the best single arm on the formal part, 2.55 on merit and 2.21 on discursive questions. None of the 180 conditions was trained on a mixed corpus, so this value is a target and not a result.
\subsection{Block composition matters more than block size}
The blocks do not form a ladder of volume, and this changes what the contrasts between them measure. The contrast \textbf{I against II} is not one of volume: the two blocks are disjoint and of similar size, 1,007 against 1,182, so it compares \emph{replacing} official Bar material with generated material. The result is weak on every task, with a mean gain of -0.03 points on the formal part and one significant arm out of nine; swapping the official for the generated buys nothing. The contrast \textbf{II against III} adds 525 cascaded synthetic documents and is a genuine addition: null on discursive questions, with two significant arms out of nine and a mean gain of +0.02, and significant in nine out of nine on the formal part, with +4.44 on average. The contrast \textbf{I against IV} adds all the generated material to the official, 1,769 examples, and has the largest effect, with eight arms out of nine on the formal part and +5.90 mean gain, against one arm out of nine on discursive questions and +0.31. Finally \textbf{IV against VI}, which adds the 223 examples from other career examinations, is null on every task.
None of the four contrasts isolates volume from composition, because each addition brings material of a kind that was not there. What holds is that adding generated material to the official one pays off, and pays off on document writing; that replacing the official with the generated does not; and that the last addition, of material from another genre, pays off nowhere. The spread across blocks says the same: the distance between the best and worst column of each row runs from 4.3 to 6.2 points on discursive questions and from 18.5 to 28.8 on the formal part. What the blocks teach differently is, above all, document form.
Block V is the sharpest case. It holds 223 examples from legal career examinations and no Bar material, and in Tables 2 and 3 it is the worst column of nearly every row, with the deficit concentrated on the formal part. The explanation is genre, and it makes a testable prediction: a judgment has no addressing, because a judge addresses no one, and neither does a legal opinion. Block V lowers the addressing criteria by 19.4 to 26.0 points across all four models, without exception.
\subsection{Comparison with the base model}
The tables give the experiment's internal starting point. It remains to compare them with the model as published, without pre-training and without fine-tuning. On \textbf{document writing} the pipeline works on all four models, and the gain concentrates on the formal cut: from 17.5 to 22.9 percentage points, against 5.6 to 11.8 on merit. On \textbf{discursive questions} it works on the two Tucano2 models, with 3.9 and 6.8 points, and fails on the two Qwen3 models: on the 8B only one of the twelve trained conditions beats the base, and on the 14B none does, the best falling 4.0 points short. The reading is not one of degradation but of \textbf{reallocation}. The whole pipeline teaches the form of the procedural document, and what it charges is the ability to answer a short question, an ability the Qwen3 models already had at a high level, arriving at 29.5\% and 35.5\% on discursive questions, above the best trained Tucano2.
\subsection{By area of law}
Table 4 fixes Tucano2 3.7B and opens the seven areas. The pattern of Section 4.1 holds in all of them: both pre-training regimes pay off on document writing and neither pays off on discursive questions. One note is due on a collection defect described in the datasheet: the Civil Code entered the corpus with 57 of its 2,046 articles and the ordinary Criminal Code did not enter at all, whereas the other five areas have their fundamental statute with coverage between 89\% and 100\%. If the value of pre-training came from reading statutory text, the gain should be smaller in those two areas. It is not: Criminal Law has the largest gain of all areas on discursive questions. What holds is the negative claim, and not an assertion about the mechanism.
\section{Discussion}
The finding that organises the others is that the effect of each decision depends on the textual genre of the task, and depends to the point of reversing sign. This has a direct methodological consequence for the domain-adaptation literature: reporting an aggregate score over a heterogeneous evaluation set can produce the conclusion that an intervention does not work when it works well on one genre and poorly on another. In our case the II-against-III contrast is null on discursive questions, with two significant arms out of nine, and significant in nine out of nine on the formal part of document writing.
The decomposition is worth one step beyond genre. Within document writing, form and merit respond differently to the same intervention, and the aggregate score, which is 85.1\% merit, records only one of the two. Whenever the evaluated task divides into subtasks of a different nature, and an examining board's rubric is a ready-made division, the aggregate score is a weighted mean that reflects the more numerous subtask and not necessarily the one of interest.
The second finding is that the form of the corpus is a design variable, not an implementation detail. Restructuring material into cards is neither uniformly better nor uniformly worse than using the text as it is: each form teaches what its material contains. The correlation of $\rho$ = 0.30 between the gains and the better than two-to-one ratio between exclusive and shared criteria indicate that choosing between the two is the wrong question, and that what matters is the combination, which was not trained here.
The third has a practical implication for those building pipelines. Continued pre-training on raw legal text costs response-formatting ability before it yields anything, and only pays for itself if enough fine-tuning follows. Reports that do not measure the intermediate state may be attributing to pre-training a gain that is, in large part, the recovery capacity of the following stage.
As for scale, the difference tracks prior exposure to Portuguese more than parameter count. The Portuguese-focused models gain on both tasks; the multilingual ones gain a great deal on document writing and lose on discursive questions. Additional material is new information for a model that lacked it and competition for capacity in one that had it.
\section{Conclusion}
We measured separately three decisions that the current recipe for domain adaptation treats as one, in a factorial design of 180 conditions evaluated against an official itemised rubric. We conclude the following. Scores from distinct textual genres should not be aggregated in evaluations of domain adaptation, and the same caution applies to subtasks of a different nature within a genre. Raw text and restructured material are complements, not substitutes: they improve sets of criteria that overlap little. And continued pre-training, in the domain and task measured, delivers predominantly document form, charging beforehand a degradation that only the subsequent fine-tuning recovers. We report the following limitations. The control regime was run across all six blocks only on Tucano2 3.7B; on the other three models it covers the extremes of volume, blocks I and VI. The two corpora are not matched on token count, 59.43 against 50.58 million, and expansion inverts the composition by category, so part of the effect attributed to form may be mixture. And the structured corpus, generated by a teacher model without reference verification, contains fabricated case citations that reach the trained model's output at a rate of 0.512\%, against 0.072\% in the control; in absolute terms it is low, in direction it is systematic, and the correction belongs to the next version of the corpus. The next step, already enabled by the material built here, is to train the combined condition, with fixed-proportion mixing, two-stage training and genre routing as the three forms to test, and to verify how much of the ceiling measured here converts into actual gain.
\section*{Acknowledgements and declarations}
\textbf{Use of generative AI.} Generative models are objects of study here in two roles: they generated the structured corpus (Section 3.1) and graded every response (Section 3.4). The authors additionally used a language model to help write the analysis code and to draft and translate the text, composed first in Portuguese. Every figure is produced by the deposited code from the deposited judgments, and every reference was checked against its primary source. The authors take full responsibility for the content.

\begin{thebibliography}{99}
\bibitem{r1} Gururangan, S., Marasović, A., Swayamdipta, S., Lo, K., Beltagy, I., Downey, D., Smith, N. A. \emph{Don't Stop Pretraining: Adapt Language Models to Domains and Tasks}. ACL, 2020.
\bibitem{r2} Chalkidis, I., Fergadiotis, M., Malakasiotis, P., Aletras, N., Androutsopoulos, I. \emph{LEGAL-BERT: The Muppets straight out of Law School}. Findings of EMNLP, 2020.
\bibitem{r3} Colombo, P., Pires, T. P., Boudiaf, M. et al. \emph{SaulLM-7B: A pioneering Large Language Model for Law}. arXiv:2403.03883, 2024.
\bibitem{r4} Pires, R., Malaquias Junior, R., Nogueira, R. \emph{Automatic Legal Writing Evaluation of LLMs}. ICAIL '25: Proceedings of the Twentieth International Conference on Artificial Intelligence and Law, 420--424, 2025. doi:10.1145/3769126.3769227.
\bibitem{r5} Pires, R., Sales Almeida, T., Larcher Junior, C., Bonás, G., Abonizio, H., Piau, M., Malaquias Junior, R., Laitz, T., Nogueira, R. \emph{Magis-Bench: Evaluating LLMs on Magistrate-Level Legal Tasks}. arXiv:2605.08437, 2026.
\bibitem{r6} Pacheco, E. C. B., Suriani, F. M., Ribeiro, R. \emph{Rabula: A Benchmark for Evaluating LLMs in Brazilian Legal Tasks}. Proceedings of the First Argument Mining and Empirical Legal Research Workshop (AMELR 2025), CEUR Workshop Proceedings, vol. 4089, paper 6, 2025.
\bibitem{r7} Zheng, L., Chiang, W.-L., Sheng, Y. et al. \emph{Judging LLM-as-a-Judge with MT-Bench and Chatbot Arena}. NeurIPS, 2023.
\bibitem{r8} Pacheco, E., Russo, C. M. \emph{Measuring the competence of LLMs in Brazilian law}. Companion paper, in preparation, 2026.
\bibitem{r9} Hu, E. J., Shen, Y., Wallis, P. et al. \emph{LoRA: Low-Rank Adaptation of Large Language Models}. ICLR, 2022.
\bibitem{r10} Benjamini, Y., Hochberg, Y. \emph{Controlling the False Discovery Rate}. Journal of the Royal Statistical Society B, 57(1), 289--300, 1995.
\bibitem{r11} Pacheco, E., Russo, C. M. \emph{Paired raw and structured corpora for Brazilian legal domain adaptation}, v1.0 [Data set]. Zenodo, 2026. doi:10.5281/zenodo.22739216.
\bibitem{r12} Pacheco, E., Russo, C. M. \emph{Judgments of 180 trained conditions on a Brazilian Bar Examination benchmark}, v1.0 [Data set]. Zenodo, 2026. doi:10.5281/zenodo.22739184.
\bibitem{r13} Pacheco, E., Russo, C. M. \emph{Agreement of 31 candidate automatic judges against a human gold standard on Brazilian legal tasks}, v1.0 [Data set]. Zenodo, 2026. doi:10.5281/zenodo.22739161.
\bibitem{r14} Pacheco, E. \emph{Synthetic corpora in the adaptation of language models to Brazilian law} [Source code]. github.com/eduardocbpacheco/synthetic-corpora-brazilian-law, 2026.
\end{thebibliography}
\end{document}
